\section{Problemstellung}
\label{sec:problemstellung}

\subsection{Fachliche Betrachtung}
\label{subsec:fachliche-betrachtung}

Beim Rollenspiel \Abk{DSA} handelt es sich um ein \sog \Abk{PPRPG}.
Diese zeichnen sich \insb dadurch aus, dass die Teilnehmer physisch am selben Ort sind und vor allem mit Stift und Papier gearbeitet wird.
Grundsätzlich wird in einem \Abk{PPRPG} eine fiktive Geschichte von fiktiven Charakteren durch die Teilnehmer interaktiv konstruiert und erzählt.
Ein Teilnehmer besetzt die Rolle des \sog Spielleiters\; die Restlichen sind \sog Spieler.
Der Spielleiter führt anhand einer fertig beschriebenen Spielwelt und einem passenden Regelwerk die Spieler durch die Geschichte.
Die Spieler besetzen jeweils die Rolle eines Charakters und beschreiben diesen auf einen \sog Charakterbogen, auf dem \uA dessen Beruf und Fähigkeiten stehen.

Das \Abk{PPRPG} Regelwerk zu \Abk{DSA} ist im Jahr 1984 erschienen und befindet sich aktuell in der 5ten-Ausgabe aus 2015.
Es zeichnet sich \insb durch den Umfang und die Komplexität aus, was \uA die Character-Erstellung und -Verwaltung für die Teilnehmer deutlich schwer macht.

\subsubsection{Charakter-Erstellung \textit{(vereinfachte Darstellung)}}
\label{subsubsec:charakter-erstellung}

Ein Character definiert sich durch die folgenden Merkmale: seine Spezies, Kultur, Profession, Eigenschaften, Fertigkeiten sowie Vor- und Nachteilen.
Einige Merkmale können Stufen \bzw eine Wertigkeit haben, \uA Eigenschaften, Fertigkeiten sowie Vor- und Nachteile.
Fertigkeiten bestehen aus Talenten, Zaubern, magischen Handlungen und Liturgien.
Jeder dieser Merkmale, oder die Steigerung einer Stufe/Wertigkeit, kostet dem Spieler \Abk{AP} und kann bei der Wahl von Nachteilen \Abk{AP} zurückbekommen.
Bei der Character-Erstellung hat jeder Spieler eine, vom Spielleiter, vorgegebene Anzahl an Start-\Abk{AP}.
Zunächst wird damit die Spezies, Kultur(en) und Profession ausgewählt.
Mit den restlichen \Abk{AP} wird für Fertigkeiten, Vorteile und andere Spielmechaniken bezahlt oder durch Nachteile erworben.
Abschließen sind acht Basiswerte, nach vorgebenden Formeln, zu berechnen.
Diese werden \ggfs durch die Spezies, den Eigenschafts-Stufen, Vor- und Nachteilen sowie Sonderfertigkeiten beeinflusst.
\Vgl{RWiki-Erstellung}

Dabei ist aus folgenden und teils durch Regeln eingeschränkten Umfängen zu wählen (dies stellt ein Auszug und keine komplette Aufstellung dar):
\begin{itemize*}[label={}, itemjoin={,}]
    \item vier Spezien
    \item 32 Kulturen
    \item 51 Professionen (bis zu zwei)
    \item acht Eigenschaften
    \item 59 Talente
    \item 251 Zaubern und 74 Ritualen
    \item 181 magischen Handlungen
    \item 74 Liturgien
    \item 130 Vorteile
    \item 77 Nachteile
    \item und weiteren Spielmechaniken wie \zB 21 Kampftechniken und 354 Sonderfertigkeiten.
\end{itemize*}
\Vgl[ff.]{RWiki}

Die Komplexität entsteht durch die Vielzahl der Merkmals-Kombinationen und Regeln.
Alle Regeln, ohne weitere Hilfsmittel, einzuhalten ist aufgrund der Anzahl für die Teilnehmer eine zeitintensive und fehleranfälligen Aufgabe.

Die \og Merkmale können sich \bspw gegenseitig beeinflussen oder ausschließen:
\begin{itemize}
    \item Eine Spezies gibt Boni und Mali für Eigenschaften vor \Vgl{RWiki-Spezies-Elfen}.
    \item Spezies und Professionen können übliche und unübliche Kulturen sowie Vor- und Nachteile vorgeben. \Vgl{RWiki-Spezies-Elfen, RWiki-Beruf-BruderDesFeuers}
    \item Professionen und Sonderfertigkeiten haben Voraussetzungen \uA \ggfs an Eigenschaften, Fertigkeiten, Vor-und Nachteilen und Sonderfertigkeiten \Vgl{RWiki-Beruf-BruderDesFeuers,RWiki-Sonderf-Blutrune}.
    \item Eine Profession gibt Fertigkeiten (Talente, Zauber, Rituale, Liturgien) vor \Vgl{RWiki-Beruf-BruderDesFeuers}.
    \item Einige Fertigkeiten stehen nur bestimmten Professionen zur Verfügung (\sog Berufsgeheimnis*se) \Vgl{RWiki-Berufsgeheimnis}.
\end{itemize}

Weitere Regeln, welche bei der Charakter-Erstellung zu beachtet sind, werden nachfolgend (und nicht
abschließend) aufgezeigt:
\begin{itemize}
    \item Es darf eine vorher festgelegte maximale Anzahl an aktivierten Zauber und Liturgien, sowie \textquote{Fremdzauber} (nicht der Profession zugeordnet) nicht überschritten werden. \Vgl[Schritt 2]{RWiki-Erstellung}
    \item Eigenschaften, Fertigkeiten und Kampftechniken können nur bis zu einem vorab festgelegten maximalen Wert verbessert werden. \Vgl[Schritt 2]{RWiki-Erstellung}
    \item Die erste Kultur ist gratis. \Vgl[Schritt 4]{RWiki-Erstellung}
    \item Zauber und Liturgien sind vor der Wert-Steigerung zu aktivieren, dafür sind Aktivierungskosten zu zahlen. \Vgl[Schritt 8]{RWiki-Erstellung}
    \item Es darf nur bis zu einer bestimmten Anzahl \Abk{AP} Vorteile gekauft oder mit Nachteilen erworben werden. \Vgl{RWiki-VorUndNachteile}
    \item Am Ende darf nur eine maximale Anzahl \Abk{AP} übrig bleiben. \Vgl[Schritt 11]{RWiki-Erstellung}
\end{itemize}

Bei den letzteren Zwei handelt es sich um Vorgaben der Spielwelt, damit die Merkmals-Kombinationen in dieser stimmig sind.
Unübliche Merkmals-Kombinationen sind nicht per se verboten, müssen jedoch bekannt und mit dem Spielleiter abgesprochen werden.

Die nachstehende \Formel{eq:merkmal-kombinationen} verdeutlicht vereinfacht und beispielhaft die Vielzahl der Merkmals-Kombinationen.
Es werden die \og Umfänge für Merkmale genutzt und komplexe Regeln, wie \zB mögliche gegenseitige Merkmals-Ausschlüsse, nicht berücksichtigt.
Gemäß Regelwerk für einen Spieler mit dem Erfahrungsgrad \textquote{Erfahren} gilt, dass jede Eigenschaft um bis zu sechs Stufen gesteigert werden kann und insgesamt 36 Eigenschaftsstufen erkauft werden können \Vgl[Schritt 2 und 5]{RWiki-Erstellung}.
Es wird angenommen, dass eine zweit Profession gewählt werden kann, sechs Eigenschaften maximal gesteigert werden, 20 Fertigkeiten aktiviert und gesteigert werden und 10 Vorteile und fünf Nachteile gewählt werden.
%
%Gemäß Regelwerk für einen Spieler mit dem Erfahrungsgrad \textquote{Erfahren} gilt, dass es 1.100 Start-\Abk{AP} gibt, jede
%Eigenschaft auf Stufe 8 startet, die Stufen-Obergrenze je Eigenschaft bei 14 liegt (also um bis zu 6 Stufen gesteigert
%werden kann), jede Eigenschaftsstufe 15 \Abk{AP} kostet, die maximale Eigenschaftsstufen-Summe bei 100 liegt, jede
%Fertigkeit maximal auf Stufe 10 (Stufe 0 nach Aktivierung) erweitert werden kann, maximal je 80 \Abk{AP} an Vorteile
%gekauft oder Nachteile gewonnen werden darf. \ZitatVgl[Schritt 2, 5 und 7]{RWiki-Erstellung}\\
%Angenommen wird das jeder Vorteil und Nachteil einen \Abk{AP}-Wert von 8 hat, 10 Vorteile und 5 Nachteile gewählt
%werden und 300 \Abk{AP} für die freie Aktivierung und Steigerung von Fertigkeiten zur Verfügung steht, wodurch 100
%Fertigkeits-Aktivierungen und Steigerungen vorgenommen werden.
%
% 8x 8 Stufen = 64 Stufen von 100 ==> also maximal 36 neue Stufen
% 8x : 9 10 11 12 13 14 = 6 Stufen * 15 AP = 90 AP
% Ich darf 8 mal 0 bis 6 Stufensteigerungen waehlen, maximal aber nur 36 Stufensteigerungen
% 36 Stufen aus 8 gruppen von bis 6 Stufen => Multinomialkoeffizient
%
% xx [F]ertigkeiten, jede kann um bis zu 11 Stufen (incl. Aktivierung) gesteigert werden, 100 stehen zur verfügung
% xx + 100 -1   über  100
% (F * 11)! + 100 - 1   über  100
%
%% TODO halbiere als annahme vorteile und nachteile, sowie fertigkeiten, wegen auschlussregeln
\begin{equation}
    \label{eq:merkmal-kombinationen}
    \underset{[S]pezien}{
        \binom{S}{1}
    }  *
    \underset{[K]ulturen}{
        \binom{K}{1}
    }  *
    \underset{[P]rofessionenn}{
        \left( P + \binom{P}{2} \right)
    }  *
    \underset{[E]igenschaften}{
    %% TODO mach was draus
        \left( \binom{36}{6, 6, 6, 6, 6, 6} \frac{36!}{6!^{E}} \right)
    } *
    \underset{[F]ertigkeiten}{
        \binom{F}{20}
    } *
    \underset{[V]orteile}{
        \binom{V}{10}
    } *
    \underset{[N]achteile}{
        \binom{N}{5}
    }
\end{equation}
%% TODO korrigiere ergebnis entsprechend
Daraus ergeben sich \textbf{\ca 2,0684 Quattuordezilliarden (10\textsuperscript{87}) verschiedene Merkmals-Kombinationen}.
Unter Fertigkeiten wurden folgende Spielmechaniken zusammengefasst: Talente, Zauber, Rituale, magischen Handlungen, Liturgien, Kampftechniken und Sonderfertigkeiten.

\subsubsection{Charakter-Verwaltung \textit{(vereinfachte Darstellung)}}
\label{subsubsec:charakter-verwaltung}

Bei der Charakter-Verwaltung geht es um Zukäufe und Steigerungen von Merkmalen, welche mit, im Laufe der Spielgeschichte, erhaltenen \Abk{AP} getätigt werden können.
Die meisten der bereits, zur Charakter-Erstellung, genannten Regeln und einhegenden Problemen gelten auch bei der Verwaltung.
So \zB Merkmals-Ausschlüsse, bei dem ein bereits gewähltes/gekauftes Merkmal ein anderen nicht erlaubt (sinngemäß nach unüblichen Merkmalen).
Oder das Einhalten der Regeln, ohne Hilfsmittel, eine zeitintensive und fehleranfällige Aufgabe ist.
Zusätzlich zu den bekannten kauf- \bzw steigerbaren Merkmalen können drei der acht Basiswerte gesteigert werden.
%% TODO fill me further
\Vgl{RWiki-Verwaltung}

\subsection{Technische Betrachtung}
\label{subsec:technische-betrachtung}

Um den zeitlichen Aufwand zu reduzieren, veröffentlichte bereits im Jahr 1993 und 1995 der damalige Herausgeber eine Software mit dem Namen \textquote{\Abk{DSA}-Tools} zur Erstellung und Verwaltung von Charakteren für die damalige Ausgabe \Vgl{Wiki-DSA-Tools,Wiki-DSA-Tools-Deluxe}.
Die aktuell genutzten Softwares kommen hingegen aus der Community selbst und stehen kostenlos zur Verfügung, nachfolgend die populärsten aufgelistet:
\begin{itemize}
    \item \textbf{Optolith Heldengenerator}, erschienen am \Datum{2018-04-08} und zuletzt Aktualisiert am \Datum{2022-03-22}, dient der Charakter-Erstellung und -Verwaltung in der 5ten-\Abk{DSA}-Ausgabe.
    Es ist seit \Monat{2020-09} OpenSource und wird hauptsächlich von einer Person gewartet.
    Genutzt werden die Sprachen TypeScript\Note{Typisierter Superset von JavaScript; \url{https://www.typescriptlang.org/}} und ReasonML\Note{OCaml (\url{https://ocaml.org/}) in JavaScript-angelehnter Notation; \url{https://reasonml.github.io/}}.
    Die Regeln werden in \AbkK{YAML}-Dateien festgehalten und ein Zugriff wird, aufgrund von Lizenzgebungen, nur auf Anfrage und Notwendigkeit ermöglicht.
    \Vgl{GitHub-Optolith}

    \item \textbf{Helden-Software}, erschienen im Jahr \Jahr{2012} und zuletzt Aktualisiert am \Datum{2017-07-23}, dient der Charakter-Erstellung und -Verwaltung in der 4ten-\Abk{DSA}-Ausgabe (für 5te-Ausgabe ist keine Unterstützung geplant \Vgl{HS-DSA5}).
    Es ist nicht OpenSource.
    Genutzt wird die Sprache Java\Note{\Abk{OS}-unabhängige \Abk{OOP} Sprache; \url{https://www.oracle.com/de/java/}} und keine weiteren Frameworks/Bibliotheken. \Vgl{HS-QA-DevEnv}

    Es wird von einem Team aus sieben Personen (fünf Entwickler und zwei Tester) gewartet \Zitat{HS-Team}.
    Als Besonderheit bietet diese Lösung das Speichern der Charakter auf deren Server und eine Schnittstelle zur Entwicklung von Plugins an, mit denen Informationen zu den Charakteren abgerufen werden können \Vgl{HS-Internet,HS-Plugins}.

    \item \textbf{DSA Excel Heldenblatt}, erschienen im Jahr \Jahr{2003} und zuletzt Aktualisiert am \Datum{2021-05-15}, dient der Charakter-Erstellung und -Verwaltung in der \Abk{DSA}-4.1 Ausgabe.
    Es handelt sich hierbei um eine Tabellenkalkulation, basierend auf Excel\Note{Vom Hersteller Microsoft; \url{https://www.microsoft.com/de-de/microsoft-365/excel}}, weshalb dieses als OpenSource zu betrachten ist.
    Die Regelwerk-Werte hingehen sind, wie auch bei den bereits genannten Programmen, aufgrund von Lizenzen verschlüsselt und damit nicht frei zugänglich.
    Das Besondere hier ist, dass eigene Merkmals-Ausprägungen wie \zB Zauber, Rituale und Liturgien integriert werden können.
    \Vgl{Orkenspalter-Heldenblatt,Heldenblatt}

    \item \textbf{DSA MeisterGeister}, erschienen im Jahr \Jahr{2010} und zuletzt Aktualisiert am \Datum{2022-06-06}, dient dem Spielleiter als Unterstützung in verschiedenen Bereich, darunter auch der Charakter-Erstellung und -Verwaltung, in der \Abk{DSA}-4.1 Ausgabe (für die 5te-Ausgabe ist keine Unterstützung geplant \Vgl{MG-DSA5}) \Vgl{MG}.
    Erst mit der letzten Aktualisierung wurde der Programmcode veröffentlicht, aufgrund fehlender personellen Kapazitäten im Team, von aktuell zwei Personen \Vgl{MG-OpenSource, MG-Team}.
    Es ist in C\#\Note{Typsichere \Abk{OOP} Sprache; \url{https://learn.microsoft.com/de-de/dotnet/csharp/tour-of-csharp/}} geschrieben und daher auf das \Abk{OS} \textquote{Microsoft Windows} ausgelegt \Vgl{MG-GitHub}.
\end{itemize}

%% TODO zeige was im generellen die technische schwierrigkeit ist
%% TODO zeige was schwierigkeit bei opensoruce ist (hindernisse von contributoren)

%%%  Arten von Rollenspiele:
%%%  https://de.wikipedia.org/wiki/Rollenspiel_(Spiel)
%%%
%%%  Defenition PP:
%%%  https://de.wikipedia.org/wiki/Pen-%26-Paper-Rollenspiel
%%%
%%%  PP wissenschaftlich:
%%%  https://de.wikipedia.org/wiki/Rollenspieltheorie

